%==============================================================
\section{正則型との比較：平坦なシンプレクティックファイバー束}\label{sec:regular}
%==============================================================
これまでの対象（型が跳ぶ）と対照的な、型が一定の「正則型」を扱い、次元の増え方が質的に違うことを示す。

\subsection{設定と $\omega$ への非依存}
\begin{defn}[設定 (S)]\label{defn:S}
$X$ をコンパクト複素多様体、$\tilde X\to X$ を普遍被覆、$\Gamma=\pi_1(X)$、$(Y,\omega_Y)$ をコンパクト連結なシンプレクティック多様体とする。準同型 $\rho:\Gamma\to\mathrm{Symp}(Y,\omega_Y)$ に対し $M=(Y\times\tilde X)/\Gamma$ とおく。「$Y$ 上のシンプレクティック構造と $X$ 上の複素構造の積」の一般化複素構造は $\Gamma$-不変なので $M$ に降り、型 $k=\dim_\C X$ の正則型の一般化複素構造が得られる。葉はファイバー $Y_x$、葉空間は $X$。$E=\pi^*G$（$G$ は $X$ 上の正則直線束）を SGH（強一般化正則）直線束とする。\tC\ 正則型・SGH の定義：Pal--Poddar \cite{PP} Def.\ 3.1 ほか（番号は\tQ）。
\end{defn}
この設定には、直積 $Y\times X$（$\Gamma$ 自明）、$S^2$-束（ホロノミーが $SO(3)$ に入るもの）、Kodaira--Thurston 型の $T^2$-束などが含まれる。$\HH(E)=H^*(M;C(L;E))$ は第 2 節と同様に定義される。

\begin{prop}\tP\label{prop:omega}
$\tilde\omega$ を $\mathrm{pr}_Y^*\omega_Y$ が降りて得られる $M$ 上の閉 2-形式とする。$e^{i\tilde\omega}$ は Courant 自己同型で、$L=e^{i\tilde\omega}L_0$（$L_0=T_{F,\C}\oplus T^{0,1}_{\mathrm{hor}}\oplus N^*$、$N^*=\pi^*T^{*1,0}_X$）。よって $L\cong L_0$ は Lie 亜群として同型で、一般化 Dolbeault 複体は
\[C(L;E)=\Omega^\bullet_{F_\C}(M;E\otimes\wedge^\bullet T^{1,0}_X),\qquad F_\C=T_{F,\C}\oplus T^{0,1}_{\mathrm{hor}}\ (\text{包合的})\]
と同一視でき、$N$ は SGH 束に自明に作用する。コホモロジーは $\omega$ を見ない（葉付き複素構造だけを見る）。
\end{prop}
\begin{proof}
局所的には積の主張で、すべての対象が $\Gamma$-不変。$d\tilde\omega=0$ なので $e^{i\tilde\omega}$ は Courant 括弧を保つ。$N$ が自明に作用するのは SGH の定義（横断方向に正則、葉に沿って平坦）による。\qedhere
\end{proof}

\subsection{主定理：Leray と Lefschetz}
\begin{thm}[正則型の $\HH^*$]\tP\ \tQ\label{thm:L}
設定 (S) を仮定し、さらに (HL)：$(Y,\omega_Y)$ は hard Lefschetz をみたす（例えば $Y$ が Kähler）とする。$V^j$ を $H^j(Y)$ を係数とする平坦束（Gauss--Manin）、$E=\pi^*G$ とすると
\[\HH^k(M,E)\ \cong\ \bigoplus_{a+b+j=k}H^a\bigl(X,\ \wedge^bT_X\otimes V^j\otimes G\bigr).\]
$X$ に Kähler や射影性は要らない。
\end{thm}
\begin{hyptab}
(S) 平坦なファイバー束 & 命題 \ref{prop:omega} による分解 $L\cong L_0$、Ehresmann 的な局所自明性（ファイバーのコホモロジーが局所系 $V^j$ になる） & 手順 (i)\\\midrule
$Y$ コンパクト連結 & $H^j(Y)$ が有限次元で、葉方向コホモロジー層が $V^j$ の滑らかな切断層になる & 手順 (i)\\\midrule
(HL) hard Lefschetz & $\Lambda=\omega_F\wedge$ が $V^{m_Y-i}\to V^{m_Y+i}$ の同型（平坦束の同型）を与える & 手順 (iii)：Deligne の補題で $d_r=0\ (r\ge2)$\\\midrule
$X$ コンパクト & 各 $H^a(X,\cdot)$ が有限次元 & 次元の和\\\midrule
$E=\pi^*G$（SGH） & $N$ が自明に作用するので $\wedge^bT_X$-次数 $b$ が直和因子（Poisson 的な $d_1$ がない） & 手順 (i)：$E_2$ の形
\end{hyptab}
\begin{proof}
(i)（二重複体）命題 \ref{prop:omega} により $C(L;E)=\Omega^\bullet_{F_\C}(M;E\otimes\wedge^\bullet T_X)$。これは「葉方向の de Rham」と「$X$ 上の Dolbeault（係数 $\wedge^bT_X\otimes G$）」の二重複体を作る。ファイバーがコンパクトなので葉方向コホモロジー層は $V^j$ の $C^\infty$ 切断の層であり、その上の $\bar\partial$ は Gauss--Manin 接続の $(0,1)$-部分で、持ち上げの取り方によらない（持ち上げの違いは葉方向 Lie 微分で、閉形式に対しては完全形式）。よって
\[E_2^{a,j}=\bigoplus_bH^a(X,\wedge^bT_X\otimes V^j\otimes G).\]
(ii) $\omega_F:=\tilde\omega|_F$ の混合成分は 0（$\tilde\omega$ は局所的に $\mathrm{pr}_Y^*\omega_Y$）。したがって $\Lambda=\omega_F\wedge$ は二重複体の鎖準同型（$d_F$ と $\bar\partial$ に可換）で、スペクトル系列の各項に作用し $d_r$ と可換。$E_2$ 上では各点のファイバー上の Lefschetz 作用素 $V^j\to V^{j+2}$ である。
(iii)（Deligne の補題）hard Lefschetz により $V^{m_Y-i}\to V^{m_Y+i}$ が同型。$x$ を重み $s$ の原始元とすると $\Lambda^{s+1}d_2x=0$ だが、$d_2x$ はファイバー次数 $m_Y-s-1$ の側にあり、そこでは $\Lambda^{s+1}$ は単射。よって $d_2=0$。帰納的に $d_r=0\ (r\ge2)$、次元は足し合わさる。\qedhere
\src{Deligne \cite{Del68}（Publ.\ Math.\ IH\'ES 35）。Leray スペクトル系列：Bredon \cite{Bre}（該当箇所は\tQ）。}
\end{proof}
\tQ\ 概略にとどまる点：(i) で、$\bar\partial$ が Gauss--Manin 接続の $(0,1)$-部分であることと、葉方向 Dolbeault 型の二重複体の $E_2$ の同定を詳細には書いていない。(HL) を外すと $d_2\ne0$ になりうる\tK（具体例は計算していない）。

\begin{cor}[次数別 Riemann--Roch]\tP\label{cor:L}
定理 \ref{thm:L} の仮定に加え、$X$ が射影的で $G$ が豊富とする。ある $m_0$ があって、$m\ge m_0$ ならば $H^a(X,\wedge^bT_X\otimes V^j\otimes G^m)=0\ (a>0)$、したがって
\begin{gather*}\HH^k(M,E^m)=\bigoplus_{b+j=k}H^0\bigl(X,\wedge^bT_X\otimes V^j\otimes G^m\bigr),\\ \dim\HH^k(M,E^m)=\sum_{b+j=k}b_j(Y)\int_X\ch(\wedge^bT_X)\,e^{m\,c_1(G)}\,\Td(X).\end{gather*}
右辺は各 $k$ について $m$ の多項式である。
\end{cor}
\begin{hyptab}
$G$ 豊富、$X$ 射影的 & Serre 消滅を連接層 $\wedge^bT_X\otimes V^j$ に使う & 全次数 $a>0$ の項の消滅\\\midrule
$V^j$ が平坦束 & $\ch(V^j)=\mathrm{rk}\,V^j=b_j(Y)$ & $b_j(Y)$ が外に出る\\\midrule
定理 \ref{thm:L} & $\HH^k$ を $\bigoplus H^a$ に分解 & 最初の等式\\\midrule
HRR & $a>0$ が消えるので $H^0$ の次元が $\chi$ に等しい & 最後の等式
\end{hyptab}
\begin{proof}
Serre 消滅 \tC\ \cite{Hart}（III.5.2）を局所自由層 $\wedge^bT_X\otimes V^j$ に適用し、有限個の $(b,j)$ の最大の $m_0$ を取る。残りは HRR と $\ch(V^j)=b_j(Y)$。\qedhere
\end{proof}
全 Euler 標数は $\sum_k(-1)^k\dim\HH^k=\chi(Y)\cdot(-1)^{\dim X}\chi(X)$ で、定理 \ref{thm:E}（のこの設定への拡張）と整合する\tV（例で確認）。各次数ごとの次元は $m$ の多項式として増えるが、全体の交代和は $m$ によらない。

\subsection{Pal--Poddar の定理との関係}\label{subsec:PPcompare}
\begin{thm}[Pal--Poddar, Thm.\ 11.12]\tC
$M$ をコンパクトな正則型一般化複素多様体、$M/\mathcal S$ を Kähler 軌道体、$E$ を正の SGH 直線束とする。(1) $p+q>k$ ならば $H^q(M,(\mathcal G^*M)^p\otimes E)=0$。(2) 任意の SGH 束 $E'$ に対し、$m\ge m_0$、$q>0$ なら $H^q(M,E'\otimes E^m)=0$。
\src{\cite{PP}, Thm.\ 11.12（仮定・主張は \texttt{pp\_recheck.md} での再読による）。}
\end{thm}
\begin{prop}\tP\ \tQ\label{prop:PPcompare}
設定 (S) では、Pal--Poddar の群 $H^q(M,(\mathcal G^*M)^p\otimes E)$ は $H^q(X,\Omega^p_X\otimes G)$（葉に沿って定数な横断形式）として読める。そのとき (1) は $X$ 上の古典的な Kodaira--Nakano 消滅、(2) は $E'=\pi^*(\wedge^bT_X\otimes V^j)$（SGH 束）に対する古典の Serre 消滅であり、系 \ref{cor:L} の消滅 $H^a(X,\wedge^bT_X\otimes V^j\otimes G^m)=0\ (a>0)$ と同じである。したがって横断次数での消滅は Pal--Poddar Thm.\ 11.12(2) を $E'=\wedge^bT_X\otimes V^j$ に適用したものに一致し、系 \ref{cor:L} は「Pal--Poddar の Serre 型消滅 $+$ Leray $+$ Lefschetz」である。
\end{prop}
\begin{rem}[\tQ]
文字通りの $M$ 上の層コホモロジーを取ると、$R^j\pi_*\pi^{-1}\ne0$ のために葉方向因子 $H^j(Y)$ が余分に現れる。Pal--Poddar の Thm.\ 8.9, Cor.\ 8.10 は $H^q(M,(\mathcal G^*M)^p)$ を横断 Dolbeault コホモロジーと同定するので、彼らの $H^q$ は葉空間上で読むべきと思われるが、原文でそのように明記されているかは確認していない。また $\HH^k$ の消滅が言えるのは横断次数 $a>0$ についてだけで、$H^j_{dR}(Y)$ 因子は各次数 $j$ に残る。
\end{rem}

\subsection{例}
\begin{ex}[直積 $Y\times\mathbb P^1$]\tP\ \tV\label{ex:prodY}
$X=\mathbb P^1$、$G=\mathcal O(m)$、$m\ge-1$：定理 \ref{thm:L} と $H^0(\mathbb P^1,\mathcal O(m))=m+1$、$H^0(\mathbb P^1,T\otimes\mathcal O(m))=m+3$ から $P_{\HH}(t)=P_Y(t)\bigl((m+1)+(m+3)t\bigr)$。$Y=S^2$：$(m+1)(1+t^2)+(m+3)(t+t^3)$（$m=0$：$1,3,1,3$）。$Y=T^2$：$m=0$ で $1,5,7,3$、一般に $(m+1,\ 3m+5,\ 3m+7,\ m+3)$。どちらも $m$ について線形に増える。
\end{ex}
\begin{ex}[Kodaira--Thurston：$T^2\to M\to E$]\tP\ \tV\ \tQ
モノドロミーがせん断 $\psi=\begin{pmatrix}1&1\\0&1\end{pmatrix}$ のとき $V^0=V^2=\mathcal O$、$V^1=F_2$（Atiyah \cite{Atiyah} の非自明な冪単な平坦束、$h^0=h^1=1$\tC）。$G=\mathcal O$ では $P_{\HH}(t)=(1+t)^2(1+t+t^2)=1+3t+4t^2+3t^3+t^4$。独立な確認：Heisenberg 代数 $\mathfrak h_3$ の葉方向複体に $\wedge N$ を掛けた Lie 環複体の Betti 数が $(1,3,4,3,1)$ で一致\tV。この確認は不変形式のモデルが $\HH^*$ を計算するという Nomizu 型の定理を仮定しており、その定理は証明していない\tQ（定理 \ref{thm:L} 自体はそれを使わない）。$\deg G=d>0$ では $P_{\HH}(t)=d(1+t)^3$、$\deg G=0$、$G\not\cong\mathcal O$ では $P_{\HH}=0$。
\end{ex}
\begin{ex}[Bieberbach 商、$S^2$-束]\tP\ \tV
$(T^2\times E)/\mathbb Z_2$（$(y,x)\mapsto(-y,x+\tau/2)$）：$V^1=L_{-1}\oplus L_{-1}$（非自明な符号表現）は $G=\mathcal O$ で非輪状、$P_{\HH}(t)=(1+t)^2(1+t^2)$（Künneth の $\mathbb Z_2$-不変部分と一致）。ホロノミーが $SO(3)$ の $S^2$-束では $V^j$ が自明束なので、数値は直積と同じ。
\end{ex}
\begin{rem}\tP
葉方向に非自明な平坦なホロノミー指標 $\chi$ をもつ SGH 束は、コンパクトな底では指標が局所定数になる（最大値原理）。$Y$ がトーラスで $\chi\ne1$ なら $H^*(Y;L_\chi)=0$ で $\HH=0$。
\end{rem}
\tC\ これらの例の数値（Kodaira--Thurston, Bieberbach）は総合版の記録からの転記で、本論文の執筆中には再計算していない\tQ。

%--------------------------------------------------------------
\subsection{型が跳ぶ場合との対比：多項式と周期的}\label{subsec:contrast}
%--------------------------------------------------------------
\begin{center}\small
\begin{tabular}{>{\raggedright\arraybackslash}p{3.4cm}>{\raggedright\arraybackslash}p{4.4cm}>{\raggedright\arraybackslash}p{5.8cm}}\toprule
 & 正則型（$Y=S^2$, $X=\mathbb P^1$） & 型が跳ぶ（$\mathbb{CP}^2$、滑らかな 3 次曲線）\\\midrule
$m$ が大きいとき $\dim\HH^k$ & $m$ の線形関数 $(m+1),(m+3),(m+1),(m+3)$ & 有界：$m\in3\Z_{\ge0}$ で $(1,0,2,0,0)$、それ以外で $(0,0,3,0,0)$\\\midrule
$m=0$ & $(1,3,1,3)$ & $(1,0,2,0,0)$\\\midrule
$m=1$ & $(2,4,2,4)$ & $(0,0,3,0,0)$\\\midrule
$m=3$ & $(4,6,4,6)$ & $(1,0,2,0,0)$\\\midrule
全 Euler 標数 & $\chi(Y)\cdot(-1)^{k}\chi(X)=-4$（$k=\dim X=1$） & $\chi(\mathbb{CP}^2)=3$（$m$ によらない）\\\bottomrule
\end{tabular}
\end{center}
\tV\ 右列は例 \ref{ex:cp2}（$m=-15,\dots,15$ で直接計算と一致）、左列は系 \ref{cor:L}。
\begin{obs}\tP\label{obs:contrast}
両者はともに $E_1^{a,b}=H^a(\wedge^bT\otimes E)$ から出発する。正則型では SGH 束で $N$ が自明に作用するので $d_1=0$（定理 \ref{thm:L}）。型が跳ぶ場合は $d_1=[\sigma,\cdot]+V\wedge$ が $H^0$ の項の間で\emph{ほぼ完全}になり、$H^0(\wedge^bT\otimes L^m)$ が $m$ について増えるにもかかわらず、残るコホモロジーは $U=X\setminus D$ の位相で決まる有界な量になる。「$m$ の多項式」か「周期的」かの違いは、$d_1$ の消滅か（ほぼ）完全性かに帰着する。
\end{obs}
\begin{conj}[統一予想（不正確）]\tK\label{conj:U}
有限個の型の層をもつコンパクトな一般化複素多様体と、正則な層に沿って $N$ が自明に作用する加群 $E$ について、$\HH^*(E)$ は「葉空間上の葉方向 de Rham $\otimes$ 横断 Dolbeault」の層ごとの寄与に濾過され、各層で $(c)$ 型（コンパクト台）か非 $(c)$ 型かは指数根（留数）で決まる。
\end{conj}
\noindent この予想は正確な形では定式化されておらず、上の 2 族（正則型と対数シンプレクティック）以外での検証はない。対応の目安は次の表である。
\begin{center}\small
\begin{tabular}{>{\raggedright\arraybackslash}p{3.0cm}>{\raggedright\arraybackslash}p{5.4cm}>{\raggedright\arraybackslash}p{5.6cm}}\toprule
 & 正則型 & 型が跳ぶ（対数シンプレクティック）\\\midrule
葉方向 de Rham & コンパクトな葉 $Y$ の $H^j(Y)$（局所系 $V^j$） & 開葉 $U$ の $H^j(U;L_U)$\\\midrule
横断方向 & $X$ 上の Dolbeault & $D$ の各層上の複体 $Q_T$\\\midrule
境界の扱い & なし（$Y$ はコンパクト） & 留数による $Rj_*$ / $j_!$ の選択（Deligne）\\\midrule
モノドロミー & Gauss--Manin & $e^{2\pi\sqrt{-1}r}$（$D$ のまわり）\\\midrule
退化 & Lefschetz で $d_r=0$ & $U$ は非コンパクトで類似物なし、$d_1\ne0$\\\bottomrule
\end{tabular}
\end{center}
